\documentclass[11pt]{article}
\usepackage[margin=1in]{geometry}
\usepackage{amsmath,amssymb,amsthm}
\usepackage{bm}
\usepackage{xcolor}

\newtheorem{proposition}{Proposition}
\theoremstyle{remark}
\newtheorem{remark}{Remark}

\newcommand{\R}{\mathbb{R}}
\newcommand{\sm}{\operatorname{softmax}}
\newcommand{\ve}[1]{\bm{#1}}

\title{Hand-Constructed Feed-Forward Weights in the Two-Cell Induction Transformer:\\
The \emph{Identity} and \emph{Log} Readouts, and Why the Log Readout Coincides with the $L_1$ Head}
\author{Markov-noise ICL experiments (folder \texttt{8. train with mlp})}
\date{\today}

\begin{document}
\maketitle

\begin{abstract}
We give the explicit weight settings that make the position-wise feed-forward (FFN)
sub-blocks of the two-cell induction transformer implement (i)~the \emph{identity} map
(used for the middle FFN) and (ii)~a \emph{log} map (used for the last FFN, immediately
before the un-embedding and softmax). We then prove that the log readout makes the
softmax head emit exactly the divide-by-sum (``$L_1$'') distribution over the
in-context count vector. This is the reason the hand-constructed log-FFN curve and the
$L_1$ curve overlap \emph{exactly} in the comparison plot: the log precisely inverts the
softmax exponential, reducing the softmax head to linear normalization.
\end{abstract}

\section{Setup and notation}

The network embeds a length-$T$ sequence ($T=200$) over a vocabulary of size $V=20$
(signal tokens $0,\dots,V_{\mathrm{sig}}-1$ with $V_{\mathrm{sig}}=10$; the rest noise).
Token and position one-hots are concatenated into the initial residual stream of width
$D_0=V+T=220$. Each of the two \emph{cells} applies a concatenating attention block
followed by a residual FFN sub-block; widths grow as
\[
D_0=220 \;\xrightarrow{\ \text{cell 1}\ }\; D_1=D_0+d=440
\;\xrightarrow{\ \text{cell 2}\ }\; D_2=D_1+d=660 ,
\qquad d=220\ \text{(attention output width)} ,
\]
and the FFN hidden width is $m=128$. A linear un-embedding
$W_o\in\R^{V\times D_2}$ (with bias $b_o$) and a softmax produce the predictive
distribution.

\paragraph{The induction circuit produces a count vector.}
The attention weights are fixed to the hand-built perfect-induction circuit $W^\star$.
Its only property we use here is the following: at every position $t$, a designated
block of the cell-2 residual stream,
\[
\ve c \;=\; \big(s_{j_1},\dots,s_{j_V}\big)\in\R^{V},\qquad J=\{j_1,\dots,j_V\}\subset\{1,\dots,D_2\},
\]
holds the \emph{in-context empirical next-token distribution} (the attention-weighted
successor frequencies of the current state). Because it is a convex combination of
one-hot vectors,
\begin{equation}\label{eq:c}
\ve c \ge \ve 0, \qquad \textstyle\sum_{i=1}^{V} c_i \le 1 .
\end{equation}
This $\ve c$ is exactly the quantity the analytic \texttt{ctx-count} predictor uses.

\paragraph{Residual FFN sub-block.}
For a residual stream $\ve s\in\R^{D}$ the FFN sub-block computes
\begin{equation}\label{eq:ffn}
\mathrm{FF}(\ve s) \;=\; \ve s \;+\; W_2\,\phi\!\big(W_1\ve s + \ve b_1\big) + \ve b_2 ,
\qquad W_1\in\R^{m\times D},\ W_2\in\R^{D\times m},
\end{equation}
with a point-wise nonlinearity $\phi$. We write $\ve e_k$ for the $k$-th standard basis vector.

\paragraph{Two readout heads.}
Given a logit vector $\ve\ell\in\R^V$, the model uses the \emph{softmax} head
$\sm(\ve\ell)_i = e^{\ell_i}/\sum_j e^{\ell_j}$. We will compare against the
\emph{divide-by-sum} (``$L_1$'') head with per-class floor $\ve\beta\ge\ve0$,
\begin{equation}\label{eq:l1}
p^{L_1}_i \;=\; \frac{\max(z_i,0)+\beta_i}{\sum_{j}\big(\max(z_j,0)+\beta_j\big)} ,
\end{equation}
which is the head used by the \texttt{l1} construction; here $\ve z$ is the (non-negative)
count readout, i.e.\ $\ve z=\ve c$.

\section{Identity construction (the middle FFN)}

We want cell-1's FFN to leave the stream untouched, $\mathrm{FF}(\ve s)=\ve s$.

\begin{proposition}[Identity]\label{prop:id}
Setting the output projection to zero,
\[
W_2 = 0,\qquad \ve b_2 = \ve 0 ,
\]
makes \eqref{eq:ffn} the identity map, $\mathrm{FF}(\ve s)=\ve s$, for \emph{any}
$W_1,\ve b_1,\phi$.
\end{proposition}
\begin{proof}
With $W_2=0$ and $\ve b_2=\ve 0$ the second term of \eqref{eq:ffn} vanishes identically,
leaving $\mathrm{FF}(\ve s)=\ve s$.
\end{proof}

\noindent
Thus the middle FFN contributes nothing to the residual stream: the two-cell
construction collapses onto the single-FFN construction below. (Operationally one may
also zero $W_1,\ve b_1$; zeroing the \emph{output} projection is the cleanest because it
is independent of $\phi$.) Note this is a \emph{choice of construction}; a \emph{trained}
middle FFN is generally not identity and does useful nonlinear work.

\section{Log construction (the last FFN)}\label{sec:log}

We now set cell-2's FFN and the un-embedding so that the logit of token $i$ is
\begin{equation}\label{eq:target}
\ell_i \;=\; \log\!\big(c_i+\beta_i\big),
\end{equation}
where $\ve\beta=(\beta_1,\dots,\beta_V)$ is a per-class floor (a single value
$\beta_{\mathrm{sig}}$ on signal coordinates and $\beta_{\mathrm{noise}}$ on noise
coordinates). The floor is supplied by the FFN \emph{bias} and prevents the
$\log 0$ blow-up on never-observed tokens (the zero-frequency problem).

Let $F=\{f_1,\dots,f_V\}\subset\{1,\dots,D_2\}$ be a block of \emph{free} coordinates of
the cell-2 stream that are disjoint from the count block $J$ and are zero under
$W^\star$ (the cell-2 attention block has $d=220$ coordinates of which only $V=20$ carry
$\ve c$, so $200$ free coordinates are available). Set:
\begin{align}
\text{(select counts)}&& (W_1)_{i,:} &= \ve e_{j_i}^{\!\top}, & (b_1)_i &= \beta_i, & i&=1,\dots,V,\label{eq:W1}\\
\text{(nonlinearity)}&& \phi &= \log, &&&& \label{eq:phi}\\
\text{(write to free block)}&& (W_2)_{f_i,\,i} &= 1, & \ve b_2 &= \ve 0, & i&=1,\dots,V,\label{eq:W2}\\
\text{(read free block)}&& (W_o)_{i,\,f_i} &= 1, & \ve b_o &= \ve 0, & i&=1,\dots,V,\label{eq:Wo}
\end{align}
with all unspecified entries equal to zero (hidden units $V{+}1,\dots,m$ are unused).

\begin{proposition}[Log readout]\label{prop:log}
With \eqref{eq:W1}--\eqref{eq:Wo}, the model's logits satisfy
$\ell_i=\log(c_i+\beta_i)$, hence \eqref{eq:target} holds.
\end{proposition}
\begin{proof}
By \eqref{eq:W1} the pre-activation of hidden unit $i\le V$ is
$u_i=(W_1\ve s)_i+(b_1)_i=s_{j_i}+\beta_i=c_i+\beta_i$. By \eqref{eq:phi},
$\phi(u_i)=\log(c_i+\beta_i)$. By \eqref{eq:W2} the residual increment
$W_2\phi(\cdot)$ is $\log(c_i+\beta_i)$ at coordinate $f_i$ and $0$ elsewhere; since the
free coordinates start at $0$, the updated stream obeys
$\tilde s_{f_i}=0+\log(c_i+\beta_i)$. Finally \eqref{eq:Wo} reads exactly those
coordinates, so $\ell_i=\tilde s_{f_i}=\log(c_i+\beta_i)$. (Routing the output to the
free block $F$ rather than back onto $J$ avoids the residual term $+\,\ve c$; equivalently
one may target $J$ and write $\log(c_i+\beta_i)-c_i$, undoing the skip connection.)
\end{proof}

\section{Why the log readout coincides with the $L_1$ head}

\begin{proposition}[Exact equivalence]\label{prop:equiv}
Under Proposition~\ref{prop:log} and \eqref{eq:c}, the softmax of the constructed logits
equals the $L_1$ head \eqref{eq:l1} with the same floor $\ve\beta$:
\[
\sm(\ve\ell)_i
\;=\; \frac{c_i+\beta_i}{\sum_{j=1}^{V}\big(c_j+\beta_j\big)}
\;=\; p^{L_1}_i .
\]
\end{proposition}
\begin{proof}
Substituting $\ell_i=\log(c_i+\beta_i)$ and using $e^{\log x}=x$,
\[
\sm(\ve\ell)_i
=\frac{e^{\ell_i}}{\sum_j e^{\ell_j}}
=\frac{e^{\log(c_i+\beta_i)}}{\sum_j e^{\log(c_j+\beta_j)}}
=\frac{c_i+\beta_i}{\sum_j (c_j+\beta_j)} .
\]
For the $L_1$ head, the count readout is $\ve z=\ve c\ge\ve0$ by \eqref{eq:c}, so
$\max(z_i,0)=c_i$ and \eqref{eq:l1} gives the same expression. The two distributions
are therefore identical.
\end{proof}

\begin{remark}[The log cancels the softmax exponential]
Softmax is ``exponentiate, then normalize''. The natural log is the inverse of the
exponential, so $\sm\circ\log$ is just \emph{normalization}:
\[
\sm\!\big(\log \ve u\big)_i=\frac{u_i}{\sum_j u_j}\qquad(\ve u>\ve0).
\]
Inserting a log-FFN before the softmax therefore \emph{removes} the softmax
nonlinearity, turning the softmax head into the linear divide-by-sum ($L_1$) head. With
identical floors $\ve\beta$ the two constructions are the \emph{same function} of the
context, so they produce identical predictive distributions and hence identical
cross-entropy at \emph{every} threshold $\alpha$. This is exactly why the
hand-constructed log-FFN curve and the \texttt{l1} curve lie on top of one another in the
comparison plot.
\end{remark}

\begin{remark}[Floors $=$ smoothing]
The floor $\ve\beta$ (the FFN bias $\ve b_1$ in \eqref{eq:W1}) is the additive smoothing
of the empirical counts. It is what avoids assigning probability $\approx 0$ to
unobserved tokens; without it, $\log 0=-\infty$ makes the readout catastrophically
over-confident. Choosing $\beta_{\mathrm{sig}},\beta_{\mathrm{noise}}$ (e.g.\ by a small
grid search) makes the resulting distribution match the smoothed count predictor
\texttt{ctx-count} to within $\approx 0.01$--$0.03$ nats.
\end{remark}

\begin{remark}[Realizability of $\phi=\log$]
A standard $\mathrm{ReLU}$ or $\mathrm{GELU}$ FFN cannot compute $\log$ \emph{exactly}.
However, by \eqref{eq:c} the argument lies in the compact interval
$c_i+\beta_i\in[\beta_i,\,1+\beta_i]$, bounded away from $0$, on which $\log$ is smooth.
A $\mathrm{ReLU}$ network with $m$ hidden units approximates a smooth scalar function on
a compact interval uniformly with error $O(1/m^2)$ (piecewise-linear interpolation), so
with $m=128$ the log map is realizable to high precision; the plotted curve uses
$\phi=\log$ directly (the idealized, best-case construction).
\end{remark}

\begin{remark}[Single-FFN model]
The minimal one-FFN model (attention followed by a single residual FFN before the
readout) admits the \emph{same} log construction of Section~\ref{sec:log}. Combined with
Proposition~\ref{prop:id}, the two-cell construction with middle FFN $=$ identity and last
FFN $=$ log is functionally identical to this single-FFN construction --- which is why
the ``FFN $\times 1$'' and ``FFN $\times 2$ (mid $=$ identity)'' curves coincide.
\end{remark}

\section{Summary}

\begin{itemize}
\item \textbf{Identity FFN} (middle): $W_2=0,\ \ve b_2=\ve0$. The block is the identity;
the middle FFN contributes nothing.
\item \textbf{Log FFN} (last): select the count block ($W_1$), add the floor
($\ve b_1=\ve\beta$), apply $\phi=\log$, write to a free block ($W_2$), and read it out
($W_o$); the logits become $\log(\ve c+\ve\beta)$.
\item \textbf{Consequence:} $\sm(\log(\ve c+\ve\beta))=(\ve c+\ve\beta)/\!\sum(\ve c+\ve\beta)$,
the $L_1$ head. The log exactly inverts the softmax exponential, so the log-FFN readout
\emph{is} the divide-by-sum readout --- hence the exactly overlapping curves.
\end{itemize}

\end{document}
